% =============================================================================
%  PAES Studio: presentación de repaso (plantilla Beamer 16:9)
%  Lista para Overleaf: Nuevo proyecto → Subir proyecto → compilar con pdfLaTeX.
%
%  Estética minimalista con la paleta Zinc de la app: sin degradados, sin
%  símbolos de navegación, líneas finas y alternativas con su tecla A–E.
%  Comandos:
%    \alternativa{A}{texto}       alternativa
%    \alternativaok{C}{texto}     alternativa correcta (para la diapositiva de solución)
%    \paso{n}{texto}      paso numerado del desarrollo
%    \ejebarra{Eje}{0.42}{42 %}   barra proporcional de frecuencia
%  Entornos: cajaclave, cajatrampa, lineatiempo (filas «fecha & hecho \\»).
%
%  El generador scripts/pipeline/generate_latex_summary.py reemplaza los bloques
%  delimitados por %%<NOMBRE>%% … %%</NOMBRE>%%.
% =============================================================================
\documentclass[aspectratio=169,11pt]{beamer}

\usepackage{lmodern}        % antes de fontenc/babel: no depende de las fuentes EC
\usepackage[T1]{fontenc}
\usepackage[utf8]{inputenc}
\usepackage[spanish,es-noshorthands,es-tabla]{babel}
\spanishplainpercent        % \% sin espaciado automático (evita choques en modo matemático)
\usepackage{textcomp}
\usepackage{amsmath,amssymb}
\usepackage{booktabs}
\usepackage{array}
\usepackage[most]{tcolorbox}

% ── Paleta Zinc ──────────────────────────────────────────────────────────────
\definecolor{zinc950}{HTML}{09090B}
\definecolor{zinc900}{HTML}{18181B}
\definecolor{zinc700}{HTML}{3F3F46}
\definecolor{zinc500}{HTML}{71717A}
\definecolor{zinc300}{HTML}{D4D4D8}
\definecolor{zinc200}{HTML}{E4E4E7}
\definecolor{zinc50}{HTML}{FAFAFA}
\definecolor{esmeralda}{HTML}{047857}
\definecolor{rosa}{HTML}{BE123C}

\newcommand{\sen}{\operatorname{sen}}

% ── Tema mínimo ──────────────────────────────────────────────────────────────
\usefonttheme{professionalfonts}
\setbeamertemplate{navigation symbols}{}
\setbeamercolor{background canvas}{bg=white}
\setbeamercolor{normal text}{fg=zinc900}
\setbeamercolor{structure}{fg=zinc900}
\setbeamercolor{alerted text}{fg=zinc950}
\setbeamercolor{item}{fg=zinc500}
\setbeamerfont{frametitle}{series=\bfseries,size=\large}
\setbeamertemplate{itemize items}{\raisebox{0.15ex}{\color{zinc500}\scriptsize$\bullet$}}
\setbeamertemplate{enumerate items}{\color{zinc500}\insertenumlabel.}
\setbeamersize{text margin left=1.6cm,text margin right=1.6cm}

\makeatletter
\setbeamertemplate{frametitle}{%
  \vspace*{0.9em}%
  {\usebeamerfont{frametitle}\color{zinc950}\insertframetitle}\par
  \ifx\insertframesubtitle\@empty\else{\small\color{zinc500}\insertframesubtitle}\par\fi
  \vspace{0.25em}{\color{zinc200}\rule{\textwidth}{0.4pt}}%
}
\makeatother
\setbeamertemplate{footline}{%
  \hspace*{1.6cm}{\color{zinc500}\tiny\insertshorttitle}\hfill
  {\color{zinc500}\tiny\ttfamily\insertframenumber/\inserttotalframenumber}\hspace*{1.6cm}\par\vspace*{0.7em}%
}
\setbeamertemplate{title page}{%
  \vfill
  {\color{zinc500}\small PAES Studio}\par\vspace{0.6em}
  {\usebeamerfont{title}\bfseries\Huge\color{zinc950}\inserttitle}\par\vspace{0.5em}
  {\large\color{zinc700}\insertsubtitle}\par\vspace{1.2em}
  {\color{zinc300}\rule{0.35\textwidth}{0.4pt}}\par\vspace{0.6em}
  {\small\color{zinc500}\insertdate}\par
  \vfill
}

% ── Componentes ──────────────────────────────────────────────────────────────
\newtcbox{\tecla}{on line, boxrule=0.4pt, arc=1.5pt, colback=zinc50, colframe=zinc300,
  left=3pt, right=3pt, top=1pt, bottom=1pt, fontupper=\ttfamily\footnotesize\color{zinc700}}
\newtcbox{\teclaok}{on line, boxrule=0.4pt, arc=1.5pt, colback=esmeralda, colframe=esmeralda,
  left=3pt, right=3pt, top=1pt, bottom=1pt, fontupper=\ttfamily\footnotesize\bfseries\color{white}}
\newtcolorbox{opcion}[1][]{enhanced, boxrule=0.4pt, arc=2pt, colback=white, colframe=zinc300,
  left=6pt, right=6pt, top=3pt, bottom=3pt, before skip=3pt, after skip=3pt, #1}
\newcommand{\alternativa}[2]{\begin{opcion}\tecla{#1}\hspace{0.6em}#2\end{opcion}}
\newcommand{\alternativaok}[2]{\begin{opcion}[colframe=esmeralda!60, colback=esmeralda!5]\teclaok{#1}\hspace{0.6em}#2\end{opcion}}
\newcommand{\paso}[2]{\par\noindent\makebox[1.4em][l]{\color{zinc500}\ttfamily\small #1}\parbox[t]{\dimexpr\linewidth-1.4em}{#2}\par\vspace{0.45em}}
\newtcolorbox{cajaclave}[1][]{enhanced, boxrule=0.4pt, arc=2pt, colback=zinc50, colframe=zinc300,
  borderline west={2pt}{0pt}{zinc900}, left=7pt, right=7pt, top=3pt, bottom=4pt,
  fonttitle=\bfseries\small, coltitle=zinc900, colbacktitle=zinc50, titlerule=0pt, title={#1}}
\newtcolorbox{cajatrampa}[1][Trampas frecuentes]{enhanced, boxrule=0.4pt, arc=2pt, colback=white, colframe=rosa!35,
  borderline west={2pt}{0pt}{rosa}, left=7pt, right=7pt, top=3pt, bottom=4pt,
  fonttitle=\bfseries\small, coltitle=zinc900, colbacktitle=white, titlerule=0pt, title={#1}}
\newenvironment{lineatiempo}
  {\footnotesize\renewcommand{\arraystretch}{1.15}\begin{tabular}{@{}>{\leavevmode\bfseries\color{zinc700}}p{0.14\textwidth}p{0.74\textwidth}@{}}}
  {\end{tabular}}
\newlength{\anchobarra}
\newcommand{\ejebarra}[3]{%
  \par\noindent\makebox[0.34\linewidth][l]{\small #1}%
  \setlength{\anchobarra}{0.5\linewidth}%
  \makebox[\anchobarra][l]{{\color{zinc200}\rule[0.2ex]{\anchobarra}{1.1ex}}\hspace*{-\anchobarra}{\color{zinc900}\rule[0.2ex]{#2\anchobarra}{1.1ex}}}%
  \hspace{0.8em}{\small\ttfamily #3}\par\vspace{0.5em}}

%%<META>%%
\title[Repaso PAES: Competencia Lectora]{Repaso PAES}
\subtitle{Competencia Lectora}
\date{1 de noviembre de 2026}
%%</META>%%

\begin{document}

\begin{frame}[plain]
  \titlepage
\end{frame}

%%<CONTENIDO>%%
\begin{frame}{Contenidos}

  \framesubtitle{3 ejes, en el orden del temario}

  \begin{enumerate}\small

    \item \textbf{Localizar}: Información explícita

    \item \textbf{Interpretar}: Vocabulario contextual, Tema, idea central y síntesis, Propósito y función de fragmentos, Inferencias y relaciones entre ideas, Textos literarios, Textos discontinuos e infografías

    \item \textbf{Evaluar}: Evaluar argumentos y postura, Evaluar forma, recursos y pertinencia

  \end{enumerate}

\end{frame}

\begin{frame}{Qué evalúa la prueba}

  \framesubtitle{Distribución por eje en 7 ensayos (388 preguntas)}

  \vspace{0.8em}

  \ejebarra{Interpretar}{0.809}{81 \%}

  \ejebarra{Evaluar}{0.121}{12 \%}

  \ejebarra{Localizar}{0.070}{7 \%}

\end{frame}

\section{Localizar}

\begin{frame}[plain]
  \vfill
  {\color{zinc500}\small Eje}\par\vspace{0.3em}{\Large\bfseries Localizar}\par\vspace{0.6em}
  {\color{zinc700} Encontrar información explícita y reconocer cómo se presenta.}\par\vspace{1em}
  {\small\color{zinc500} Información explícita}
  \vfill
\end{frame}

\begin{frame}[allowframebreaks]{Información explícita}
  \framesubtitle{Teoría esencial}
  \begin{itemize}\small
    \item Busca palabras clave del enunciado en el texto.
    \item La correcta es una paráfrasis fiel, no una copia literal.
    \item Una alternativa verdadera en la realidad pero ausente del texto es incorrecta.
  \end{itemize}
  \begin{cajaclave}[Cómo reconocer la correcta]\footnotesize
La alternativa correcta \textbf{reformula} lo que dice el texto sin agregar, quitar ni cambiar el alcance. Desconfía de las que generalizan («siempre», «todos») o que usan las mismas palabras del texto con otro sentido.
  \end{cajaclave}
\end{frame}

\begin{frame}{Desafío}
  \framesubtitle{Información explícita}
  {\small «Las abejas visitan cerca de un tercio de los cultivos que consumimos. Sin embargo, el uso de ciertos plaguicidas ha reducido sus colonias en varias regiones del país.»

Según el fragmento, ¿qué ha provocado la disminución de las colonias de abejas?}
  \vspace{0.4em}
  \begin{tcbraster}[raster columns=2, raster column skip=0.6em, raster row skip=0.35em]
  \alternativa{A}{La disminución de los cultivos.}
  \alternativa{B}{El uso de ciertos plaguicidas.}
  \alternativa{C}{El cambio climático.}
  \alternativa{D}{La sobreexplotación de la miel.}
  \end{tcbraster}
\end{frame}

\begin{frame}[allowframebreaks]{Solución}
  \framesubtitle{La causa está explícita en la segunda oración.}
  {\small
  \paso{1}{Ubica la oración que habla de la reducción de colonias.}
  \paso{2}{«El uso de ciertos plaguicidas ha reducido sus colonias»: es la alternativa B.}
  \paso{3}{C es plausible en la realidad, pero el fragmento no lo menciona.}
  }
  \alternativaok{B}{El uso de ciertos plaguicidas.}
  \begin{cajatrampa}\footnotesize
    \begin{itemize}
    \item Elegir una alternativa con palabras copiadas del texto pero con otro sentido.
    \item Elegir información correcta según tus conocimientos, pero que el texto no dice.
    \end{itemize}
  \end{cajatrampa}
\end{frame}

\section{Interpretar}

\begin{frame}[plain]
  \vfill
  {\color{zinc500}\small Eje}\par\vspace{0.3em}{\Large\bfseries Interpretar}\par\vspace{0.6em}
  {\color{zinc700} Comprender el sentido global y local, inferir y relacionar.}\par\vspace{1em}
  {\small\color{zinc500} Vocabulario contextual, Tema, idea central y síntesis, Propósito y función de fragmentos, Inferencias y relaciones entre ideas, Textos literarios, Textos discontinuos e infografías}
  \vfill
\end{frame}

\begin{frame}[allowframebreaks]{Vocabulario contextual}
  \framesubtitle{Teoría esencial}
  \begin{itemize}\small
    \item Importa el sentido en el texto, no el del diccionario.
    \item La correcta debe funcionar gramaticalmente en la oración.
    \item Cuidado con los sinónimos de otra acepción de la palabra.
  \end{itemize}
  \begin{cajaclave}[Prueba de sustitución]\footnotesize
1) Lee la oración y la siguiente. 2) Formula con tus palabras qué quiere decir. 3) Reemplaza por cada alternativa. 4) Elige la que conserva el sentido exacto, no solo uno parecido.
  \end{cajaclave}
\end{frame}

\begin{frame}{Desafío}
  \framesubtitle{Vocabulario contextual}
  {\small «El ministro intentó ZANJAR la discusión con un anuncio inesperado, pero los parlamentarios siguieron debatiendo hasta la medianoche.»

¿Cuál de las siguientes palabras puede reemplazar a ZANJAR sin que cambie el sentido del texto?}
  \vspace{0.4em}
  \alternativa{A}{Excavar}
  \alternativa{B}{Iniciar}
  \alternativa{C}{Terminar}
  \alternativa{D}{Postergar}
\end{frame}

\begin{frame}[allowframebreaks]{Solución}
  \framesubtitle{«Zanjar una discusión» es darla por concluida; el «pero» confirma que no lo logró.}
  {\small
  \paso{1}{Contexto: tras el intento, los parlamentarios \textbf{siguieron} debatiendo.}
  \paso{2}{El contraste indica que el ministro quería ponerle fin.}
  \paso{3}{«Terminar la discusión» conserva el sentido; «excavar» es otra acepción de «zanjar».}
  }
  \alternativaok{C}{Terminar}
  \begin{cajatrampa}\footnotesize
    \begin{itemize}
    \item Elegir el sinónimo de otra acepción («grave» como «serio» y no como «de tono bajo»).
    \item Elegir una palabra con matiz más fuerte o más débil que el original.
    \end{itemize}
  \end{cajatrampa}
\end{frame}

\begin{frame}[allowframebreaks]{Tema, idea central y síntesis}
  \framesubtitle{Teoría esencial}
  \begin{itemize}\small
    \item Tema: frase nominal. Idea central: afirmación completa.
    \item Los ejemplos y datos apoyan la idea central; no lo son.
    \item El título suele anticipar el tema.
  \end{itemize}
  \begin{cajaclave}[Criterio de amplitud]\footnotesize
La idea central no puede ser demasiado específica (un ejemplo o dato) ni demasiado general (aplicable a cualquier texto del tema). Debe cubrir todo lo que el párrafo o texto plantea.
  \end{cajaclave}
\end{frame}

\begin{frame}{Desafío}
  \framesubtitle{Tema, idea central y síntesis}
  {\small «Muchos creen que los murciélagos son plagas. Sin embargo, una sola colonia puede consumir toneladas de insectos al año, y varias especies polinizan plantas como el agave. Además, su guano se ha usado por siglos como fertilizante.»

¿Cuál es la idea central del fragmento?}
  \vspace{0.4em}
  \begin{tcbraster}[raster columns=2, raster column skip=0.6em, raster row skip=0.35em]
  \alternativa{A}{Los murciélagos polinizan el agave.}
  \alternativa{B}{Los murciélagos son beneficiosos, pese a la mala fama que tienen.}
  \alternativa{C}{El guano de murciélago es un fertilizante.}
  \alternativa{D}{Los animales suelen ser incomprendidos.}
  \end{tcbraster}
\end{frame}

\begin{frame}[allowframebreaks]{Solución}
  \framesubtitle{Cada oración aporta un beneficio que contradice la creencia inicial.}
  {\small
  \paso{1}{La primera oración plantea la creencia; «Sin embargo» introduce la postura del texto.}
  \paso{2}{Insectos, polinización y guano son ejemplos de beneficios.}
  \paso{3}{A y C son detalles; D es demasiado general.}
  }
  \alternativaok{B}{Los murciélagos son beneficiosos, pese a la mala fama que tienen.}
  \begin{cajatrampa}\footnotesize
    \begin{itemize}
    \item Elegir un detalle llamativo como idea central.
    \item Elegir una afirmación tan general que serviría para cualquier texto.
    \end{itemize}
  \end{cajatrampa}
\end{frame}

\begin{frame}[allowframebreaks]{Propósito y función de fragmentos}
  \framesubtitle{Teoría esencial}
  \begin{itemize}\small
    \item Propósito = intención, no contenido.
    \item Mira el lugar del fragmento en el texto (inicio, desarrollo, cierre).
    \item Revisa los conectores que lo rodean.
  \end{itemize}
  \begin{cajaclave}[Verbos de propósito]\footnotesize
Ilustrar o ejemplificar, introducir, contrastar, concluir, justificar, refutar, describir, advertir. La alternativa correcta combina el verbo adecuado con el contenido correcto.
  \end{cajaclave}
\end{frame}

\begin{frame}{Desafío}
  \framesubtitle{Propósito y función de fragmentos}
  {\small «Los hábitos pequeños transforman la salud. Pedro, por ejemplo, comenzó a caminar diez minutos al día y, al cabo de un año, había bajado su presión arterial.»

¿Qué función cumple la mención de Pedro?}
  \vspace{0.4em}
  \alternativa{A}{Introducir un tema nuevo.}
  \alternativa{B}{Contradecir la idea anterior.}
  \alternativa{C}{Ejemplificar la afirmación previa.}
  \alternativa{D}{Concluir el texto.}
\end{frame}

\begin{frame}[allowframebreaks]{Solución}
  \framesubtitle{El conector «por ejemplo» anuncia la función.}
  {\small
  \paso{1}{La primera oración afirma que los hábitos pequeños transforman la salud.}
  \paso{2}{El caso de Pedro muestra esa afirmación en una situación concreta: la ejemplifica.}
  }
  \alternativaok{C}{Ejemplificar la afirmación previa.}
  \begin{cajatrampa}\footnotesize
    \begin{itemize}
    \item Elegir una alternativa que resume el fragmento en vez de explicar para qué está.
    \item Acertar el verbo pero no el contenido (ilustrar, sí, pero otra idea).
    \end{itemize}
  \end{cajatrampa}
\end{frame}

\begin{frame}[allowframebreaks]{Inferencias y relaciones entre ideas}
  \framesubtitle{Teoría esencial}
  \begin{itemize}\small
    \item La inferencia correcta se apoya en el texto, no en suposiciones.
    \item Descarta alternativas que exageran o generalizan.
    \item El conector revela la relación lógica.
  \end{itemize}
  \begin{cajaclave}[Conectores frecuentes]\footnotesize
Causa: porque, ya que, debido a. Consecuencia: por lo tanto, así, en consecuencia. Contraste: sin embargo, pero, no obstante, en cambio. Adición: además, asimismo.
  \end{cajaclave}
\end{frame}

\begin{frame}{Desafío}
  \framesubtitle{Inferencias y relaciones entre ideas}
  {\small «Desde que la biblioteca amplió su horario hasta las 22:00, la sala de estudio se llena cada noche, sobre todo en época de exámenes.»

¿Qué se puede inferir del fragmento?}
  \vspace{0.4em}
  \alternativa{A}{Antes la biblioteca cerraba antes de las 22:00.}
  \alternativa{B}{La biblioteca tiene más libros que antes.}
  \alternativa{C}{Los estudiantes prefieren estudiar de noche que de día.}
  \alternativa{D}{La biblioteca abre solo en época de exámenes.}
\end{frame}

\begin{frame}[allowframebreaks]{Solución}
  \framesubtitle{«Amplió su horario hasta las 22:00» implica que antes cerraba más temprano.}
  {\small
  \paso{1}{Ampliar significa extender: el horario anterior terminaba antes.}
  \paso{2}{C exagera: el texto no compara preferencias de día y de noche.}
  }
  \alternativaok{A}{Antes la biblioteca cerraba antes de las 22:00.}
  \begin{cajatrampa}\footnotesize
    \begin{itemize}
    \item Elegir una conclusión posible, pero no necesaria.
    \item Confundir causa con consecuencia.
    \end{itemize}
  \end{cajatrampa}
\end{frame}

\begin{frame}[allowframebreaks]{Textos literarios}
  \framesubtitle{Teoría esencial}
  \begin{itemize}\small
    \item Caracteriza a un personaje por sus acciones y palabras.
    \item El narrador en primera persona solo conoce lo que vive.
    \item Interpreta la figura según el contexto del relato.
  \end{itemize}
  \begin{cajaclave}[Figuras frecuentes]\footnotesize
Comparación: «sonó como una respiración». Metáfora: «la ciudad es un hormiguero». Personificación: «el viento susurraba». Hipérbole: «te lo dije mil veces».
  \end{cajaclave}
\end{frame}

\begin{frame}{Desafío}
  \framesubtitle{Textos literarios}
  {\small «Mi abuelo era un roble: había resistido dos guerras, tres mudanzas y la muerte de su mejor amigo sin que nadie lo viera llorar.»

¿Qué se destaca del abuelo con la expresión «era un roble»?}
  \vspace{0.4em}
  \alternativa{A}{Su gran estatura.}
  \alternativa{B}{Su fortaleza ante la adversidad.}
  \alternativa{C}{Su amor por la naturaleza.}
  \alternativa{D}{Su edad avanzada.}
\end{frame}

\begin{frame}[allowframebreaks]{Solución}
  \framesubtitle{La metáfora se explica en lo que sigue: resistió hechos duros sin quebrarse.}
  {\small
  \paso{1}{Identifica la metáfora: el abuelo es un roble (árbol resistente).}
  \paso{2}{Los ejemplos que siguen (guerras, mudanzas, pérdidas) apuntan a la fortaleza.}
  }
  \alternativaok{B}{Su fortaleza ante la adversidad.}
  \begin{cajatrampa}\footnotesize
    \begin{itemize}
    \item Atribuir al narrador ideas de un personaje.
    \item Interpretar literalmente una expresión figurada.
    \end{itemize}
  \end{cajatrampa}
\end{frame}

\begin{frame}[allowframebreaks]{Textos discontinuos e infografías}
  \framesubtitle{Teoría esencial}
  \begin{itemize}\small
    \item El título dice qué representa el gráfico.
    \item Fíjate en las unidades y en la fuente.
    \item Integra la información de dos bloques cuando la pregunta lo pide.
  \end{itemize}
  \begin{cajaclave}[Orden de lectura]\footnotesize
Título $\rightarrow$ subtítulos y leyendas $\rightarrow$ datos destacados $\rightarrow$ notas y fuentes. Pregunta: ¿qué relación muestra el gráfico o la imagen con el texto?
  \end{cajaclave}
\end{frame}

\begin{frame}{Desafío}
  \framesubtitle{Textos discontinuos e infografías}
  {\small Una infografía titulada «Agua que no ves» indica: «1 kg de carne de vacuno: 15 000 L de agua. 1 kg de trigo: 1500 L. 1 taza de café: 140 L».

¿Qué propósito tiene presentar estas cifras?}
  \vspace{0.4em}
  \alternativa{A}{Mostrar cuánta agua se usa al producir alimentos.}
  \alternativa{B}{Recomendar no consumir café.}
  \alternativa{C}{Comparar precios de alimentos.}
  \alternativa{D}{Explicar cómo se cultiva el trigo.}
\end{frame}

\begin{frame}[allowframebreaks]{Solución}
  \framesubtitle{El título («Agua que no ves») y las cifras apuntan al agua usada en la producción.}
  {\small
  \paso{1}{Las cifras son litros de agua asociados a cada alimento.}
  \paso{2}{El título destaca que esa agua no es visible para el consumidor: informa sobre el agua usada al producir.}
  }
  \alternativaok{A}{Mostrar cuánta agua se usa al producir alimentos.}
  \begin{cajatrampa}\footnotesize
    \begin{itemize}
    \item Leer un dato de la columna o fila equivocada.
    \item Ignorar las notas o la fuente, que pueden cambiar la interpretación.
    \end{itemize}
  \end{cajatrampa}
\end{frame}

\section{Evaluar}

\begin{frame}[plain]
  \vfill
  {\color{zinc500}\small Eje}\par\vspace{0.3em}{\Large\bfseries Evaluar}\par\vspace{0.6em}
  {\color{zinc700} Valorar argumentos, recursos y la relación entre forma, contenido y propósito.}\par\vspace{1em}
  {\small\color{zinc500} Evaluar argumentos y postura, Evaluar forma, recursos y pertinencia}
  \vfill
\end{frame}

\begin{frame}[allowframebreaks]{Evaluar argumentos y postura}
  \framesubtitle{Teoría esencial}
  \begin{itemize}\small
    \item Primero identifica la tesis; después evalúa cada alternativa contra ella.
    \item La concesión reconoce algo del adversario antes de refutarlo.
    \item «Debilita» exige que el dato afecte directamente el argumento central.
  \end{itemize}
  \begin{cajaclave}[Estrategias argumentativas]\footnotesize
Ejemplificación, cita de autoridad, datos o estadísticas, analogía, concesión («es cierto que\ldots{} pero\ldots{}»), contraargumentación y apelación a valores.
  \end{cajaclave}
\end{frame}

\begin{frame}{Desafío}
  \framesubtitle{Evaluar argumentos y postura}
  {\small «Las ciudades deberían reemplazar los estacionamientos del centro por ciclovías: así se reduce la congestión, porque más personas dejarían el auto.»

¿Cuál de las siguientes afirmaciones, de ser verdadera, DEBILITA el argumento?}
  \vspace{0.4em}
  \alternativa{A}{Las ciclovías son más baratas que los estacionamientos.}
  \alternativa{B}{En ciudades que lo hicieron, la cantidad de autos en el centro no varió.}
  \alternativa{C}{Andar en bicicleta mejora la salud.}
  \alternativa{D}{Muchas personas ya usan bicicleta en el centro.}
\end{frame}

\begin{frame}[allowframebreaks]{Solución}
  \framesubtitle{El argumento supone que más personas dejarían el auto; B muestra que no ocurre.}
  {\small
  \paso{1}{Tesis: reemplazar estacionamientos por ciclovías reduce la congestión.}
  \paso{2}{Razón: más personas dejarían el auto.}
  \paso{3}{B contradice esa razón con evidencia; A y C hablan de otros beneficios.}
  }
  \alternativaok{B}{En ciudades que lo hicieron, la cantidad de autos en el centro no varió.}
  \begin{cajatrampa}\footnotesize
    \begin{itemize}
    \item Elegir una alternativa sobre el tema que no toca el argumento central.
    \item Confundir la opinión del emisor con la de una voz citada.
    \end{itemize}
  \end{cajatrampa}
\end{frame}

\begin{frame}[allowframebreaks]{Evaluar forma, recursos y pertinencia}
  \framesubtitle{Teoría esencial}
  \begin{itemize}\small
    \item Identifica primero el propósito y el destinatario.
    \item Juzga el recurso por su aporte, no por tu gusto.
    \item Un título adecuado anticipa el tema y la intención.
  \end{itemize}
  \begin{cajaclave}[Criterios]\footnotesize
Pertinencia: ¿aporta al propósito? Adecuación: ¿corresponde al destinatario y la situación? Coherencia: ¿es consistente con el resto del texto?
  \end{cajaclave}
\end{frame}

\begin{frame}{Desafío}
  \framesubtitle{Evaluar forma, recursos y pertinencia}
  {\small Un afiche municipal para niños sobre el lavado de manos usa dibujos coloridos y la frase «¡Jabón y agua, y los bichos se van a casa!».

¿Por qué el lenguaje del afiche es adecuado?}
  \vspace{0.4em}
  \alternativa{A}{Porque usa términos científicos precisos.}
  \alternativa{B}{Porque es simple y lúdico, acorde a su público infantil.}
  \alternativa{C}{Porque está escrito en un registro formal.}
  \alternativa{D}{Porque presenta estadísticas de enfermedades.}
\end{frame}

\begin{frame}[allowframebreaks]{Solución}
  \framesubtitle{La adecuación se juzga respecto del destinatario: niños.}
  {\small
  \paso{1}{Destinatario: niños; propósito: promover el lavado de manos.}
  \paso{2}{La rima y los «bichos» simplifican el mensaje y lo vuelven memorable.}
  }
  \alternativaok{B}{Porque es simple y lúdico, acorde a su público infantil.}
  \begin{cajatrampa}\footnotesize
    \begin{itemize}
    \item Elegir una valoración subjetiva sin respaldo en el texto.
    \item Ignorar el destinatario al evaluar el registro.
    \end{itemize}
  \end{cajatrampa}
\end{frame}
%%</CONTENIDO>%%

\end{document}
